% !TeX program = xelatex
% UTF-8 transcription of all ten source photographs, in handwritten page order.
% Source filenames (original JPEG bytes are preserved):
% pic01.jpg = 20261016_044129360.jpg
% pic02.jpg = 20261016_044136132.jpg
% pic03.jpg = 20261016_044140365.jpg
% pic04.jpg = 20261016_044145591.jpg
% pic05.jpg = 20261016_044148546.jpg
% pic06.jpg = 20261016_044154786.jpg
% pic07.jpg = 20261016_044158803.jpg
% pic08.jpg = 20261016_044204386.jpg
% pic09.jpg = 20261016_044207722.jpg
% pic10.jpg = 20261016_044217250.jpg
% Bracketed notes identify text cut off at a photograph edge; no missing words
% have been silently supplied. Medical statements/doses reproduce the notes.
% Compile standalone with XeLaTeX, or input into a XePersian document.
\ifdefined\settextfont
  \def\SeborrheicStandalone{0}
\else
  \def\SeborrheicStandalone{1}
  \documentclass[12pt,a4paper]{article}
  \usepackage[margin=2cm]{geometry}
  \usepackage{amsmath,amssymb}
  \usepackage{xepersian}
  \IfFontExistsTF{Amiri}{\settextfont{Amiri}}{\settextfont{FreeSerif}}
  \setlatintextfont{FreeSerif}
\fi
\providecommand{\en}[1]{\lr{#1}}
\newcommand{\sdpage}[1]{\par\bigskip\noindent\textbf{صفحهٔ #1 ـ موضوع: درماتیت سبورئیک ـ ۴۰}\par\noindent تاریخ: / /\par\medskip}
\newcommand{\sdcut}{\textnormal{[ادامه در لبهٔ عکس بریده شده است]}}
\ifnum\SeborrheicStandalone=1 \begin{document}\fi

\sdpage{۱}
\noindent\en{Rook's Dermatology (10th)}
\section*{درماتیت سبورئیک (\en{SD})}
\begin{itemize}
  \item \en{chronic skin disease}
  \item \en{Red or pink patches with superficial scaling}
  \item \en{Areas: Scalp, Face, Central chest, Anogenital}
  \item \en{Dandruff}:
    \begin{itemize}
      \item \en{Mildest form of SD (Pityriasis Capitis)}
      \item \en{No visible inflammation}
    \end{itemize}
  \item تشخیص \en{SD} براساس یافته‌ها بالینی صرف است.
  \item \en{SD} به چین‌های پوستی فلکسورال و ناحیهٔ ساب ممری تمایل دارد.
  \item الگوی درگیری \en{SD} معمولاً قرینه است.
  \item درگیری آن از بلوغ شروع می‌شود؛ یک واریانت آن از نوزادی.
  \item پاتوفیزیولوژی آن ناشناخته است ولی نقش مالاسزیا فورفور مطرح می‌باشد.
  \item \en{Sebo-psoriasis}:
    \begin{itemize}
      \item \en{Psoriasis-like variant of SD}
      \item \en{inflammation + well-defined, coarser scaling}
      \item \en{May Represent Koebnerization of psoriasis}
    \end{itemize}
\end{itemize}

\sdpage{۲}
\subsection*{اپیدمیولوژی \en{SD}}
\begin{itemize}
  \item متوسط ۴\% شیوع (۴٫۲ مردان، ۱٫۴ زنان).
  \item در افراد با پوست روشن شایع‌تر است.
  \item فرم نوزادان شناخته‌شده: \en{Cradle Cap}.
    \begin{itemize}
      \item بیشترین شیوع در مطالعات ۷۲\% در زیر ۳ ماه [خوانش عدد دست‌نویس نامطمئن].
      \item تا قبل از ۱ سالگی رفع می‌شود.
    \end{itemize}
  \item \en{SD} در ۳۵\% بیماران \en{HIV} در شروع بیماری رخ می‌دهد و در ۸۵\% بیماران با \en{AIDS}.
  \item در بیماران دریافت‌کنندهٔ عضو شایع‌تر است.
  \item در بیماری‌های نورولوژیک مزمن شایع‌تر است:
    \begin{itemize}
      \item \en{PD}، اسپاینال کورد اینجری،
    \end{itemize}
  \item شیوع ۳۱\% در سندرم \en{Down}.
\end{itemize}
\subsection*{ویژگی‌های بالینی \en{SD}}
\begin{itemize}
  \item ناحیهٔ درگیر می‌تواند بدون خارش یا با خارش بسیار شدید باشد.
  \item در ناحیهٔ \en{Scalp} علائم بیمار می‌تواند با شدت بیماری نامتناسب باشد.
  \item سیر بیماری مزمن و عودکننده است.
  \item \en{Facial SD}:
    \begin{itemize}
      \item چین‌ها را درگیر می‌کند: نازولابیال، پلک، گوش، ناحیهٔ گلابلا، بین دو ابرو.
    \end{itemize}
\end{itemize}

\sdpage{۳}
\begin{itemize}
  \item الگوی درگیری صورت: \en{Fine Scaling + red or pink patches}.
  \item کانال گوش معمولاً اندکی ملتهب است.
    \begin{itemize}
      \item \en{Otitis Externa} می‌تواند رخ دهد با عفونت ثانویهٔ باکتریال یا کاندیدا.
    \end{itemize}
  \item درگیری پلک: به حالت بلفاریت تظاهر.
    \begin{itemize}
      \item دبری می‌تواند وارد چشم شود و تحریک \en{Conjunctiva} ایجاد کند: \en{Red Eye}.
      \item \en{Consequent meibomian gland loss $\rightarrow$ Dry eye disease}
    \end{itemize}
  \item درگیری ناحیهٔ ژنیتال در هر دو جنس رخ می‌دهد.
  \item درگیری نواحی دیگر در مردان شایع‌تر است.
  \item در افراد با پوست تیره ممکن است \en{hypopigmentation} به همراه التهاب اندک رخ دهد.
  \item در هندی‌ها یک فرم خاص به نام سبورئیک ملانوزیس رخ می‌دهد که هایپرپیگمانته است.
\end{itemize}

\sdpage{۴}
\subsection*{\en{SD} شیرخواران}
\begin{itemize}
  \item \en{Cradle Cap; Napkin dermatitis}
  \item ظرف ۳ ماه اول رخ داده و ظرف ۱ سال اول رفع می‌شود.
  \item درگیری ابروها، ناحیهٔ پره‌نازال و نواحی فلکسورال شایع است و برای افتراق از درماتیت آتوپیک کمک‌کننده است.
  \item یک \en{DDX} مهم: \en{Psoriasiform Napkin Dermatitis}.
\end{itemize}
\begin{itemize}
  \item در بزرگسالان ممکن است همزمان با سایر درماتوزها رخ دهد.
  \item آن‌هایی که به غدد سباسه تمایل دارند:
    \begin{itemize}
      \item \en{Acne, Rosacea}، پیتریازیس ورسیکالر، \en{Pityrosporum Folliculitis}.
      \item \en{Onychomycosis, Tinea pedis}.
    \end{itemize}
  \item \en{SD} با درگیری منتشر به‌صورت نادر به \en{Exfoliative Dermatitis} به همراه اریترودرما پیشرفت می‌کند.
\end{itemize}
\subsection*{تشخیص افتراقی پوستهٔ سر برای \en{SD}}
\begin{itemize}
  \item پسوریازیس: پوستهٔ پسوریازیس \en{Well circumscribed}، ضخیم‌تر و نقره‌ای‌رنگ است.
  \item درگیری ناخن و پلاک‌های پسوریازیس وجود دارد.
\end{itemize}

\sdpage{۵}
\subsection*{تشخیص افتراقی \en{SD} با درماتیت تماسی آلرژیک}
\begin{itemize}
  \item \en{Patch Testing} در موارد درگیری آتیپیک و در موارد درماتیت پلک.
\end{itemize}
\subsection*{تشخیص افتراقی \en{SD} با \en{Darier disease} و بیماری \en{Hailey--Hailey}}
\begin{itemize}
  \item در \en{DD} ضایعات \en{Dome-shaped and papular} هستند.
  \item شک به \en{DD} و \en{HHD}: بیوپسی.
\end{itemize}
\subsection*{سایر \en{DDX}های \en{SD}}
\begin{itemize}
  \item درماتیت پری‌اورال: درگیری چین نازولابیال + \en{Scaling}.
    \begin{itemize}
      \item \en{inf. papules in a diamond-shaped distribution in periocular area}.
      \item کمک‌کننده به تشخیص.
    \end{itemize}
  \item پمفیگوس فولیاسه و پمفیگوس اریتماتوزیس.
    \begin{itemize}
      \item بیوپسی و بررسی ایمونوفلورسانس برای افتراق.
    \end{itemize}
  \item پیتریازیس روزئا: در \en{ddx} با واریانت پیتریازیفورم \en{SD}.
    \begin{itemize}
      \item \en{No herald patch + involvement beyond the torso}
    \end{itemize}
  \item \en{Early cutaneous T-cell lymphoma}.
  \item اریتراسما: شبیه درگیری \en{SD} در فلکسورهای بزرگ.
\end{itemize}

\sdpage{۶}
\subsection*{\en{DDX} در شیرخواران}
\begin{itemize}
  \item پسوریازیس، هیستیوسیتوز، کمبود زینک، آکرودرماتیت انتروپاتیک، بیماری \en{Leiner} (اریترودرما دسکواماتیروم)، کمبود بایوتینیداز.
\end{itemize}
\subsection*{\en{DDX} در کودکان قبل از بلوغ (و کودکان با پوسته‌ریزی سر)}
\begin{itemize}
  \item تینا کپیتیس (ریزش مو + شکستگی مو).
  \item پدیکولوزیس: شپش.
\end{itemize}
\subsection*{\en{Drug Induced/Exacerbated Seborrhoeic dermatitis}}
\begin{itemize}
  \item \en{Drugs containing sulphydryl groups}:
    \begin{itemize}
      \item \en{Captopril, D-penicillamine, Na-aurothiomalate}
    \end{itemize}
  \item \en{gold salts, lithium, buspirone, methyldopa, chlorpromazine, cimetidine}
  \item \en{HCV drugs: Ribavirin, Inf-$\alpha$}
  \item \en{Others}
\end{itemize}
\subsection*{طبقه‌بندی شدت \en{SD}}
\begin{itemize}
  \item سیستم طبقه‌بندی اختصاصی وجود ندارد.
  \item \en{SDASI: SD Area Severity Index}
\end{itemize}

\sdpage{۷}
\subsection*{مدیریت \en{SD}}
\begin{itemize}
  \item درمان قطعی وجود ندارد.
  \item درمان نگهدارندهٔ طولانی‌مدت لازم است.
  \item دسته داروهای مؤثر در \en{SD}:
    \begin{itemize}
      \item ضدقارچ‌ها، کراتولیتیک‌ها، ضدالتهاب‌ها / ایمونومدولاتورها، داروهای با منشأ گیاهی.
    \end{itemize}
\end{itemize}
\subsection*{ضدقارچ‌های موضعی}
\noindent\en{Foam, oint, gels, spray, shampoo}
\begin{itemize}
  \item خط اول درمان استاندارد.
  \item کتوکونازول ۲\%، سیکلوپیروکس اولامین ۱\%؛ مشابه هم.
  \item سرتاکونازول ۲\% (اندک کمتر)، بایفونازول ۱\% (بی‌اثر).
  \item مصرف طولانی‌مدت ممکن است مقاومت ایجاد کند.
  \item سایر ضدقارچ‌های موضعی:
    \begin{itemize}
      \item \en{Zinc Pyrithione}، پیروکتون اولامین، \en{Coal tar}، سلنیوم سولفاید.
    \end{itemize}
  \item ضدقارچ موضعی اثربخشی کمی در کنترل شوره دارد چراکه فراهمی زیستی و ماندگاری آن در محل کم است.
\end{itemize}
\subsection*{کراتولیتیک‌ها}
\begin{itemize}
  \item سالیسیلیک اسید، سولفور، \en{coal tar}، اوره، لاکتیک اسید، پروپیلن گلیکول.
  \item اثربخشی آن‌ها زیر سؤال است؛ خصوصاً در شیرخواران.
  \item سالیسیلیک اسید در شیرخواران ممنوع است.
\end{itemize}

\sdpage{۸}
\noindent یادداشت بالای صفحه: \en{Mild} ـ بتامتازون ۰٫۱\%، \en{OD}.
\subsection*{ضدالتهاب موضعی}
\begin{itemize}
  \item کورتیکواستروئید موضعی: مونوتراپی، در همراهی با ضدقارچ.
  \item درمان کوتاه‌مدت با کورتیکواستروئید موضعی \en{mild} توصیه می‌شود.
    \begin{itemize}
      \item اگر درمان مونوتراپی ضدقارچ به \en{SD} جواب نداد.
    \end{itemize}
  \item مهارکننده‌های موضعی کلسینورین:
    \begin{itemize}
      \item تاکرولیموس، پیمکرولیموس.
      \item در مواردی که مجبوریم کورتون را طولانی‌مدت بدهیم.
      \item اثر ظرف ۲ هفته.
      \item درمان نگهدارنده با تاکرولیموس ۰٫۱\%، ۲ بار در هفته مؤثر است.
      \item عوارض: خارش، التهاب، قرمزی.
      \item به معنای آلرژی نیست.
    \end{itemize}
  \item لیتیوم موضعی اثربخش است (لیتیوم ۸\%)؛ خصوصاً در \en{SD} صورت.
  \item سایر ضدالتهاب‌ها خصوصاً در \en{SD} صورت:
    \begin{itemize}
      \item کرم نیکوتینامید ۴\%.
      \item ژل مترونیدازول ۰٫۷۵\%؛ خصوصاً در بیماری که به‌طور همزمان آکنه / روزاسه دارد.
    \end{itemize}
\end{itemize}
\subsection*{درمان سیستمیک}
\begin{itemize}
  \item در موارد \en{Recalcitrant, widespread}.
  \item رژیم پالسی ایتراکونازول: \en{200 mg} خوراکی روزانه برای یک هفته سپس \en{200 mg} روزانه ۲ روز اول ماه.
\end{itemize}

\sdpage{۹}
\subsection*{جدول درمان استاندارد \en{SD (NICE)}}
\noindent نوع \en{SD}:
\subsubsection*{۱) اسکالپ / ریش}
\begin{itemize}
  \item خط اول: شامپو کتوکونازول ۲\% یا شامپو سلنیوم سولفید ۲٫۵\%، ۲ بار در هفته برای یک ماه سپس یک بار در هفته.
  \item خط دوم: زینک، \en{coal tar}، سالیسیلیک اسید.
  \item درمان کمکی: کراتولیتیک موضعی، روغن (زیتون / معدنی) \sdcut
  \item اگر خارش دارید: کورتیکواستروئید موضعی \en{potent} برای ۴ هفته.
\end{itemize}
\subsubsection*{۲) صورت و بدن در بزرگسالان}
\begin{itemize}
  \item خط اول:
    \begin{itemize}
      \item کرم کتوکونازول ۲\%، \en{OD/BD}.
      \item کرم کلوتریمازول ۱\%، \en{BD/TDS}.
      \item کرم اکونازول ۱\%، \en{BD}.
      \item کرم میکونازول ۲\%، \en{BD}.
    \end{itemize}
  \item خط دوم: کورتیکواستروئید موضعی \en{mild} برای ۱--۲ هفته.
  \item درمان کمکی: شامپو کتوکونازول ۲\%.
\end{itemize}
\subsubsection*{۳) شیرخواران}
\begin{itemize}
  \item خط اول: حمام روزانه و جدا کردن پوسته‌ها با شامپو که \sdcut{} و [باقی سطر بعد در لبهٔ عکس بریده شده است]
  \item خط دوم: کلوتریمازول به قیاس در بالا [خوانش عبارت «به قیاس» نامطمئن].
\end{itemize}

\sdpage{۱۰}
\begin{itemize}
  \item \en{SD} ممکن است نیازمند بررسی از نظر \en{HIV} باشد؛ اگر شرح حال مشکوک
  \item \en{SD} ممکن است جزئی از \en{IRIS} در \en{HIV} باشد.
    \begin{itemize}
      \item \en{Immune Reconstitution Inflammatory syndrome}
    \end{itemize}
  \item شیوع آکنه، روزاسه، پسوریازیس، درماتیت تماسی، فولیکولیت، پیتریازیس ورسیکالر در \en{SD} نسبت به مردم عادی بیشتر است.
  \item احتمال تشدید علائم در استرس شدید روانی بالاتر است. \en{Flair}.
  \item شیوع \en{SD} در افسردگی بالاتر است و اضطراب.
  \item در درموسکوپی: نباید \en{red dots and globules} که ویژگی پسوریازیس است را ببینیم.
  \item بیوپسی در تشخیص \en{SD} کمک‌کننده نیست.
\end{itemize}

\ifnum\SeborrheicStandalone=1
  \def\SeborrheicFinish{\end{document}}
\else
  \let\SeborrheicFinish\relax
\fi
\SeborrheicFinish
